\clearpage
\subunitheader{4.1: Toy Example: A five-step Transformer stress test}

Test how $\mu$P scaling transfers hyperparameters across extreme widths and depths.

\paragraph{Training protocol.}
Train for five updates with constant LR from the first update, no warmup,
no weight decay, and no gradient clipping. Use unscaled rotary position
embeddings (RoPE). Use the same five training batches
and 64 fixed validation sequences across parameterizations; evaluate in
microbatches and tune LR using validation loss after update five.

\paragraph{Width settings.}
Let $m=n/n_0$ and let $k$ be each matrix's fan-in. Initialize weight matrices
with independent centered Gaussian entries using the variances below.
\begin{center}
\begin{tabular}{l@{\hspace{2em}}c@{\hspace{2em}}c}
\toprule
Setting & Kaiming & $\mu$P \\
\midrule
Depth & \multicolumn{2}{c}{2} \\
Attention-head dimension & \multicolumn{2}{c}{64} \\
Reference width $n_0$ & \multicolumn{2}{c}{512} \\
Weights, activations, optimizer states & \multicolumn{2}{c}{FP32} \\
Embedding variance & \multicolumn{2}{c}{1} \\
Hidden-matrix variance & \multicolumn{2}{c}{$1/k$} \\
Initial normalization gains & \multicolumn{2}{c}{1} \\
\midrule
Readout variance & $1/n$ & $1/n_0$ \\
Readout forward multiplier & $1$ & $1/m$ \\
Hidden-matrix LR & $\eta$ & $\eta/m$ \\
Embedding, readout, norm LR & $\eta$ & $\eta$ \\
\bottomrule
\end{tabular}
\end{center}

\paragraph{Depth settings.}
Fix width and reference width at 64, with one attention head and reference
depth two. Keep parameters, gradients, Adam states, and residual additions
in FP32; use BF16 autocast for branch computations.
Let $L$ count Transformer layers, $r=L/2$, and $c_L$ multiply each attention
and feedforward branch before residual addition. Apply these depth multipliers
on top of the $\mu$P width prescription:
\begin{center}
\begin{tabular}{lccc}
\toprule
 & Residual $c_L$ & Block LR multiplier & Block Adam $\epsilon$ \\
\midrule
Standard $\mu$P & $1$ & $1$ & $10^{-8}$ \\
\href{https://arxiv.org/abs/2310.02244}{Depth-$\mu$P} & $r^{-1/2}$ & $r^{-1/2}$ & $10^{-8}r^{-1/2}$ \\
\href{https://arxiv.org/abs/2505.01618}{CompleteP} & $r^{-1}$ & $1$ & $10^{-8}r^{-1}$ \\
\bottomrule
\end{tabular}
\end{center}
Block multipliers apply to hidden matrices and within-block normalization.
Embedding, readout, and final-normalization LR and stabilizer remain unchanged.

\needspace{12\baselineskip}
\begin{problem}{Test width and depth scaling over five updates (about 55 runs)}
\paragraph{Measure feature movement, update alignment, and readout--feature-change alignment.}
Let $h_t$ be a hidden feature vector of width $n$, evaluated on $N$ fixed
probes $x_1,\ldots,x_N$. For a matrix $A_t$ with fan-in $k$, let
$\Delta A_t=A_t-A_{t-1}$ be its actual update and $u_{t-1}$ its pre-update input.
Measure feature movement from initialization and update alignment by
\[
 M_t=\left[\frac1{Nn}\sum_{i=1}^N
 \|h_t(x_i)-h_0(x_i)\|_2^2\right]^{1/2},\qquad
 \alpha_{\rm upd}^{(t)}=\log_k R(\Delta A_t,u_{t-1}).
\]
Pool feature RMS over sequences, token positions, and coordinates; compute
weight RMS over entries. Inspect normalized features and the residual stream
separately.

For readout probes, take $h_t$ to be the final RMSNorm output before the
readout multiplier, and let $V_0\in\mathbb R^{n\times q}$ be the initial
readout weights in the transposed convention. Measure
\[
 \omega_{\rm move}^{(t)}=\log_n S(V_0,h_t-h_0).
\]
This probes the initial-readout/feature-change term in the derivation.
Use the same fixed probes and pool RMS as above; omit the ratio when
feature movement is zero.

\begin{parts}
\item \textbf{Transfer across widths (about 30 five-step runs).}
Predict how well Kaiming and standard $\mu$P will transfer their best base LR
across widths 640, 2,560, and 5,120. Choose an LR sweep for each width
and parameterization.
Fit each loss--LR curve with a suitable model. Plot fitted optimal base LR
and fitted minimum validation loss against width. Which parameterization
transfers its tuned base LR better? Does it also achieve lower loss?

\item \textbf{Probe width scaling (reuse part (a)'s configurations).}
At each width and parameterization's best sampled LR, track logit RMS and
the RMS and movement $M_t$ of the final RMSNorm output on fixed probes
through the five updates. Plot update alignment $\alpha_{\rm upd}^{(t)}$ for
each block's query, key, value, attention-output, and feedforward gate/up/down
matrices, and the final readout. Use each matrix's pre-update input and its
actual fan-in $k$. Also plot $\omega_{\rm move}^{(t)}$, comparing with the no-alignment ($1/2$) and
full-alignment ($1$) scales. Compare feature movement and alignment across
layers and widths. Which assumptions and predictions of the preceding derivation agree
with your measurements?

\needspace{7\baselineskip}
\item \textbf{Transfer across depths (about 25 five-step runs).}
Predict which depth prescription should transfer its base LR most reliably.
Compare standard $\mu$P, Depth-$\mu$P, and CompleteP at depths 2, 100, and
1,000, choosing an LR sweep for each depth and prescription.
Plot fitted optimal base LR and fitted minimum validation loss against depth.
Does better LR transfer imply lower loss after five updates?

\item \textbf{Probe depth scaling (reuse part (c)'s configurations).}
At each depth and prescription's best sampled LR, track residual-stream RMS
before normalization and branch-output RMS before multiplication by $c_L$.
Measure movement of the normalized residual-stream features at matched
relative positions, such as the midpoint and final layer.
Also track $\omega_{\rm move}^{(t)}$ through
all five updates and compare it across depths and prescriptions.
Can similar normalized activations hide a growing residual stream?
Relate these probes to the LR-transfer results in part (c), distinguishing
bounded signals from nonvanishing feature changes.
\end{parts}
\end{problem}
